% !PRINTED 460
\weberpage{460}

\weq{3}{\Theta = \frac{\psi(\xi)}{F'(\xi)},}

und $\Theta$ ist also in $\Omega(\xi)$ enthalten. Da umgekehrt auch jede
Grösse in $\Omega(\xi)$ zugleich in $\Omega(\alpha, \beta, \gamma \dots)$ enthalten ist, so sind
beide Körper identisch, d.\,h. es ist

\[
\Omega(\xi) = \Omega(\alpha, \beta, \gamma \dots).
\]

Damit ist aber unsere Behauptung erwiesen.

\begin{center}
§. 144.

\gesperrt{Primitive und imprimitive Körper.}
\end{center}

Nach dem zuletzt bewiesenen Satze beschränken wir uns
jetzt auf die Betrachtung algebraischer Körper $\Omega(\alpha)$, worin $\alpha$
die Wurzel einer Gleichung in $\Omega$ ist. Diese Gleichung möge,
von mehrfachen Factoren befreit, mit

\weq{1}{F(x) = 0}

bezeichnet und vom Grade $m$ sein. Ihre Wurzeln seien

\weq{2}{\alpha, \alpha_1, \alpha_2 \dots \alpha_{m-1}.}

Es kann $F(x)$ in $\Omega$ reducibel sein; dann wird es einen und
nur einen \gesperrt{irreducibeln Factor} $f(x)$ von $F(x)$ geben, so
dass $\alpha$ eine Wurzel der Gleichung

\weq{3}{f(x) = 0}

ist. Der Grad von $f(x)$ sei $n$ und die Wurzeln von (3), die alle
unter den Wurzeln (2) enthalten sind, seien

\weq{4}{\alpha, \alpha_1, \alpha_2 \dots \alpha_{n-1}.}

Der Grad des Körpers $\Omega(\alpha)$ ist dann nach der in §.\,142
gemachten Festsetzung gleichfalls $n$.

Jede der Grössen (4) giebt zu einem algebraischen Körper
Anlass, und so entstehen die Körper

\weq{5}{\Omega(\alpha), \Omega(\alpha_1) \dots \Omega(\alpha_{n-1}),}

die wir die mit $\Omega(\alpha)$ \gesperrt{conjugirten Körper} nennen. Es kann
sein, dass diese Körper alle oder theilweise identisch sind, sie
können aber auch alle von einander verschieden sein, wovon
später Näheres.

Nach §.\,142 erhalten wir jede Grösse $\Theta$ in $\Omega(\alpha)$, wenn wir
in einer ganzen rationalen Function $\chi(t)$ in $\Omega$, höchstens vom
Grade $(n-1)$, für die Variable $t$ die Zahl $\alpha$ setzen.
